\section*{Capítulo \ref{ch:b2:solid-liquid-diagrams} --- Diagramas binários sólido--líquido}

\begin{solution}{exo:b2:solid-liquid-diagrams:1}
A \qty{1300}{\degreeCelsius}, o \omterm{def:b2:solid-liquid-diagrams:liquidus}{liquidus} está em 0.541 e o \omterm{def:b2:solid-liquid-diagrams:liquidus}{solidus} em 0.609.
0.40: só líquido. 0.58: líquido (0.541) e \omterm{def:b2:solid-liquid-diagrams:solid-solution}{solução sólida} (0.609). 0.70:
só \omterm{def:b2:solid-liquid-diagrams:solid-solution}{solução sólida}.
\end{solution}

\begin{solution}{exo:b2:solid-liquid-diagrams:2}
Fração sólida $(0.58 - 0.541)/(0.609 - 0.541) = 0.57$: \qty{1.15}{mol} de sólido,
\qty{0.85}{mol} de líquido. Níquel: $1.15 \times 0.609 = \qty{0.70}{mol}$ no
sólido, $0.85 \times 0.541 = \qty{0.46}{mol}$ no líquido; total \qty{1.16}{mol}
$= 2.00 \times 0.58$.
\end{solution}

\begin{solution}{exo:b2:solid-liquid-diagrams:3}
A pressão fixa, $v = n + 1 - \varphi$ (sem reação). (a) $1 + 1 - 2 = 0$.
(b) $2 + 1 - 2 = 1$. (c) Líquido e duas \omterm{def:b2:solid-liquid-diagrams:solid-solution}{soluções sólidas}: $2 + 1 - 3 = 0$.
\end{solution}

\begin{solution}{exo:b2:solid-liquid-diagrams:4}
Uma queda regular até \qty{232}{\degreeCelsius}, um patamar horizontal enquanto
o estanho solidifica, depois de novo uma queda regular, mais lenta perto da
temperatura ambiente. No patamar, líquido e sólido coexistem a pressão fixa,
\omterm{def:b2:equilibrium-shifts:variance}{variância} 0: o calor retirado é fornecido pela cristalização, e a temperatura
não pode mudar até que o último líquido tenha solidificado.
\end{solution}

\begin{solution}{exo:b2:solid-liquid-diagrams:5}
Benzene: $\ln x_B = -(9870/8.314)(1/269.6 - 1/278.64) = -0.1429$, $x_B =
0.867$. Naftaleno: $\ln x_N = -(19\,100/8.314)(1/269.6 - 1/353.2) = -2.017$,
$x_N = 0.133$. A soma é 1.000: \qty{269.6}{K} ($\qty{-3.5}{\degreeCelsius}$) é
a temperatura eutética, e o líquido eutético tem 0.133 de naftaleno.
\end{solution}

\begin{solution}{exo:b2:solid-liquid-diagrams:6}
O líquido esfria até o \omterm{def:b2:solid-liquid-diagrams:liquidus}{liquidus} do lado do chumbo, onde aparecem cristais de
(Pb); o líquido se enriquece em estanho, até o eutético a \qty{182.2}{\degreeCelsius}
e 62.1~\%. Ali, o líquido restante solidifica na \omterm{def:b2:solid-liquid-diagrams:eutectic}{mistura eutética}.
Fração mássica de eutético: $(30 - 18.91)/(62.13 - 18.91) = 0.26$. Com o
resfriamento posterior, a solubilidade do estanho em (Pb) diminui (linha
tracejada): um pouco de (Sn) precipita dentro dos cristais de (Pb).
\end{solution}

\begin{solution}{exo:b2:solid-liquid-diagrams:7}
23.3~\% de NaCl em massa: $23.3/76.7 = \qty{0.304}{kg}$ de sal por quilograma
de água. A \qty{-25}{\degreeCelsius}, abaixo do eutético, nenhum líquido pode
existir, quaisquer que sejam as proporções: o sal e o gelo ficam lado a lado
como sólidos.
\end{solution}

\begin{solution}{exo:b2:solid-liquid-diagrams:8}
$x_{\mathrm B} = 0.5$ fica entre E$_1$ e o composto: $\mathrm{AB_2}$
cristaliza primeiro, o líquido se desloca para E$_1$ e ali solidifica no
eutético de A e $\mathrm{AB_2}$; no final, A + $\mathrm{AB_2}$.
$x_{\mathrm B} = 0.9$ fica entre E$_2$ e B: B cristaliza primeiro, o líquido
atinge E$_2$; no final, $\mathrm{AB_2}$ + B.
\end{solution}

\begin{solution}{exo:b2:solid-liquid-diagrams:9}
Na borda traseira da zona, o líquido solidifica em um sólido que contém apenas
0.78 vez a sua fração molar de cobre: o cobre rejeitado fica no líquido, que
avança e se enriquece; assim o cobre é varrido para a extremidade da barra.
Com uma razão próxima de 1, cada passagem retira apenas uma fração da
impureza: são necessárias muitas passagens, cada uma partindo da barra
deixada pela anterior.
\end{solution}

\begin{solution}{exo:b2:solid-liquid-diagrams:10}
Dedução como no capítulo. Para um líquido diluído, $\ln x_A \approx -x_B$, $1/T -
1/T^* \approx -\Delta T/T^{*2}$, logo $x_B = \Delta_{\text{fus}}H_A\Delta T/(RT^{*2})$;
com $x_B \approx n_B/n_A = bM_A$, $\Delta T = (RT^{*2}M_A/\Delta_{\text{fus}}H_A)\,b$.
Benzeno: $K_f = 8.314 \times 278.64^2 \times 0.07811/9870 = \qty{5.11}{K.kg/mol}$.
\end{solution}

\begin{solution}{exo:b2:solid-liquid-diagrams:11}
Por \qty{100}{g}: $19.0/24.31 = \qty{0.782}{mol}$ de Mg, $81.0/207.2 =
\qty{0.391}{mol}$ de Pb; razão $2.00$: $\ce{Mg2Pb}$.
\end{solution}

\begin{solution}{exo:b2:solid-liquid-diagrams:12}
O patamar é proporcional à massa de líquido eutético: a liga desconhecida
contém $120/300 = 0.40$ dele. Do lado do chumbo, $w = 18.91 + 0.40 \times (62.13 -
18.91) = 36.2~\%$ de estanho; do lado do estanho, $w = 96.59 - 0.40 \times (96.59 - 62.13) =
82.8~\%$. A primeira quebra da \omterm{def:b2:solid-liquid-diagrams:cooling-curve}{curva de resfriamento} (mais alta para a liga de
36~\%, cujo \omterm{def:b2:solid-liquid-diagrams:liquidus}{liquidus} fica do lado íngreme do chumbo) ou uma micrografia
(cristais primários de (Pb) ou de (Sn)) permite distingui-las.
\end{solution}

\begin{solution}{pb:b2:solid-liquid-diagrams:1}
\textbf{1.} Chumbo \qty{327.5}{\degreeCelsius}, estanho \qty{231.9}{\degreeCelsius}.
\textbf{2.} Por \qty{100}{g}: $62.13/118.71 = \qty{0.5234}{mol}$ de estanho e
$37.87/207.2 = \qty{0.1828}{mol}$ de chumbo; $x_{\mathrm{Sn}} = 0.5234/0.7062 =
0.741$.
\textbf{3.} Pontos (0, 327.5), (100, 231.9), (62.13, 182.2), (18.91, 182.2),
(96.59, 182.2).
\textbf{4.} Como na figura: líquido; L + (Pb); L + (Sn); (Pb); (Sn); (Pb) + (Sn).
\textbf{5.} L (62.1~\% Sn) $\longrightarrow$ (Pb) (18.9~\%) + (Sn) (96.6~\%).
\textbf{6.} $v = 2 + 1 - 3 = 0$: a temperatura e as três composições ficam
fixadas.
\textbf{7.} A \omterm{def:b2:solid-liquid-diagrams:solid-solution}{solução sólida} rica em chumbo (Pb): o chumbo dissolve estanho;
por isso o sólido em equilíbrio com o líquido já contém estanho, na
proporção lida no \omterm{def:b2:solid-liquid-diagrams:liquidus}{solidus}.
\textbf{8.} Ela segue o \omterm{def:b2:solid-liquid-diagrams:liquidus}{liquidus} em direção ao eutético, enriquecendo-se em
estanho, até 62.13~\%.
\textbf{9.} Líquido com 62.13~\% e (Pb) com 18.91~\% de estanho.
\textbf{10.} (Pb): $(62.13 - 40)/(62.13 - 18.91) = 0.512$; líquido 0.488.
\textbf{11.} (Pb) com 18.91~\% e (Sn) com 96.59~\%; (Sn): $(40 - 18.91)/(96.59 -
18.91) = 0.271$; (Pb): 0.729.
\textbf{12.} Constituinte eutético 48.8~\% (o líquido presente logo acima);
(Pb) primário 51.2~\%. O (Pb) da questão 11 é a soma dos cristais primários
e das lamelas de (Pb) do eutético.
\textbf{13.} Uma quebra no \omterm{def:b2:solid-liquid-diagrams:liquidus}{liquidus}, uma queda mais lenta e depois um patamar a
\qty{182.2}{\degreeCelsius} que dura 0.488 vez o da liga eutética.
\textbf{14.} $\ln x_{\mathrm{Sn}} = -(7030/8.314)(1/T - 1/505.1)$.
\textbf{15.} $\ln x_{\mathrm{Bi}} = -(11\,300/8.314)(1/T - 1/544.6)$.
\textbf{16.} $x_{\mathrm{Sn}}(T) + x_{\mathrm{Bi}}(T) = 1$.
\textbf{17.} A \qty{110}{\degreeCelsius}: $0.587 + 0.349 = 0.936$; a
\qty{120}{\degreeCelsius}: $0.621 + 0.382 = 1.003$; a
\qty{130}{\degreeCelsius}: $0.655 + 0.417 = 1.072$. A soma chega a 1 perto de
\qty{119.5}{\degreeCelsius}: \qty{120}{\degreeCelsius} com precisão de um grau.
\textbf{18.} $x_{\mathrm{Bi}} = 0.38$; em massa, $0.38 \times 208.98/(0.38 \times
208.98 + 0.62 \times 118.71) = 0.52$.
\textbf{19.} O modelo dá \qty{120}{\degreeCelsius} e 52~\% de bismuto; o
diagrama avaliado, \qty{138.8}{\degreeCelsius} e 57~\%: cerca de
\qty{19}{\degreeCelsius} abaixo.
\textbf{20.} O estanho que cristaliza não é puro, mas uma solução que contém
bismuto: seu \omterm{def:b2:chemical-potential:chemical-potential}{potencial químico} é menor que o do estanho puro, o sólido é mais
estável e cristaliza a partir do líquido a uma temperatura mais alta para uma
dada composição do líquido. O ramo do estanho do \omterm{def:b2:solid-liquid-diagrams:liquidus}{liquidus} sobe e encontra o
ramo do bismuto mais acima (um líquido não ideal acrescenta sua própria
correção).
\textbf{21.} Ela funde e solidifica a uma só temperatura, a mais baixa do
sistema: não há faixa pastosa durante a qual uma junta movida trincaria.
\textbf{22.} O chumbo é tóxico, e o lixo eletrônico acaba no meio ambiente.
\textbf{23.} Componentes sensíveis ao calor podem ser soldados a uma
temperatura mais baixa; mas a junta amolece a uma temperatura mais baixa em
serviço, e o bismuto a torna quebradiça.
\textbf{24.} \qty{570}{g} de bismuto e \qty{430}{g} de estanho.
\textbf{25.} O modelo de líquido ideal prevê um eutético a $\boldsymbol{T_E
\approx \qty{120}{\degreeCelsius}}$, contra \qty{138.8}{\degreeCelsius}
avaliados.
\end{solution}
